\section{Markov Chain Monte Carlo Methods}

\subsection{Bayesian Inference}

\subsubsection*{Prior, likelihood, and posterior}

For data $x=(x_1,\ldots,x_n)$, parameters $\theta=(\theta_1,\ldots,\theta_r)$, likelihood $L(\theta)=f(x\mid\theta)$, and prior $p(\theta)$, Bayes' theorem gives

$\pi(\theta\mid x)=\dfrac{f(x\mid\theta)p(\theta)}{f(x)}\propto L(\theta)p(\theta),\qquad f(x)=\int f(x\mid\theta)p(\theta)d\theta$.

A common estimator is the posterior mean $\hat\theta=\int\theta\pi(\theta\mid x)d\theta$. A posterior can serve as the prior when new data arrive.

\subsubsection*{Normal-normal conjugacy}

If $x_i\stackrel{iid}{\sim}N(\mu,\sigma^2)$ with known $\sigma^2$ and $\mu\sim N(m,\tau^2)$, then $\mu\mid x\sim N(m_1,\tau_1^2)$, where
$m_1=\dfrac{\tau^2\bar x+m\sigma^2/n}{\tau^2+\sigma^2/n},\qquad \tau_1^2=\dfrac{\tau^2\sigma^2}{n\tau^2+\sigma^2}$.

As $n\to\infty$, $m_1\to\bar x$ and $\tau_1^2\to0$. As $\tau^2\to\infty$, the prior becomes noninformative and $\mu\mid x\sim N(\bar x,\sigma^2/n)$.

A conjugate prior produces a posterior in the same family, with updated hyperparameters. Common pairs are Poisson--Gamma, Binomial--Beta, normal mean--Normal, and normal variance--Inverse Gamma. The lecture uses

$f_G(x)=\dfrac{\beta^\alpha x^{\alpha-1}e^{-\beta x}}{\Gamma(\alpha)},\quad f_{IG}(x)=\dfrac{\beta^\alpha x^{-\alpha-1}e^{-\beta/x}}{\Gamma(\alpha)}$.

\subsubsection*{Importance sampling for posteriors}

For an easy density $q$, draw $\theta_i\sim q$ and estimate

$\hat\theta=\int\theta\dfrac{\pi(\theta\mid x)}{q(\theta)}q(\theta)d\theta\approx\dfrac1n\sum_{i=1}^n\dfrac{\theta_i\pi(\theta_i\mid x)}{q(\theta_i)}$.

Here IS facilitates integration rather than primarily reducing variance; poor overlap between $q$ and the posterior may cause large variance.

\subsection{Gibbs Sampling}

\subsubsection*{General algorithm}

MCMC constructs a Markov chain whose limiting distribution is the target posterior. Gibbs sampling cycles through full conditionals:

$\theta_j^{(i)}\sim p(\theta_j\mid\theta_1^{(i)},\ldots,\theta_{j-1}^{(i)},\theta_{j+1}^{(i-1)},\ldots,\theta_r^{(i-1)},x)$.

Under standard regularity conditions, $\theta^{(i)}$ converges in distribution to $p(\theta\mid x)$. Discard $M$ burn-in draws and estimate
$\hat\theta_j=\dfrac1n\sum_{i=1}^n\theta_j^{(M+i)}$.

\subsubsection*{Bivariate normal}

For $(\theta_1,\theta_2)^\top\sim N(0,\Sigma)$ with unit variances and correlation $\rho$,

$\theta_1\mid\theta_2\sim N(\rho\theta_2,1-\rho^2),\qquad \theta_2\mid\theta_1\sim N(\rho\theta_1,1-\rho^2)$.

Thus, using independent $Z_1,Z_2\sim N(0,1)$,

$\theta_1^{(i)}=\rho\theta_2^{(i-1)}+\sqrt{1-\rho^2}Z_1,\qquad \theta_2^{(i)}=\rho\theta_1^{(i)}+\sqrt{1-\rho^2}Z_2$.

\subsubsection*{Unknown normal mean and variance}

For $x_j\sim N(\mu,\sigma^2)$ with both unknown,

$\mu_i\mid\sigma_{i-1}^2\sim N(m_i,\tau_i^2),\quad m_i=\dfrac{\tau_{i-1}^2\bar x+m_{i-1}\sigma_{i-1}^2/n}{\tau_{i-1}^2+\sigma_{i-1}^2/n},\quad \tau_i^2=\dfrac{\tau_{i-1}^2\sigma_{i-1}^2}{n\tau_{i-1}^2+\sigma_{i-1}^2}$.

$\sigma_i^2\mid\mu_i\sim IG(\alpha_i,\beta_i),\quad \alpha_i=\alpha_{i-1}+n/2,\quad \beta_i=\beta_{i-1}+\dfrac12\sum_{j=1}^n(x_j-\mu_i)^2$.

Alternate these draws, discard burn-in, and average the retained values.

\subsection{Jump-Diffusion Case Study}

\subsubsection*{Model and latent jumps}

Merton's model is $d\log S_t=\mu dt+\sigma dW_t+YdN_t$, where $N_t$ has intensity $\lambda$ and $Y\sim N(k,s^2)$. For small $\Delta t$,

$x_i=\Delta\log S_i=\mu\Delta t+\sigma\Delta W_i+Y_i\Delta N_i,\quad P(\Delta N_i=1)=\lambda\Delta t$.

Conditional return laws are

$x_i\mid\Delta N_i=0\sim N(\mu\Delta t,\sigma^2\Delta t),\qquad x_i\mid\Delta N_i=1\sim N(\mu\Delta t+k,\sigma^2\Delta t+s^2)$.

If $f_0,f_1$ denote these densities, the latent indicator probability is

$P(\Delta N_i=1\mid x_i)=\dfrac{f_1(x_i)\lambda\Delta t}{f_1(x_i)\lambda\Delta t+f_0(x_i)(1-\lambda\Delta t)}$.

When $\Delta N_i=1$, the latent jump size satisfies

$Y_i\mid x_i\sim N\left(\dfrac{(x_i-\mu\Delta t)/(\sigma^2\Delta t)+k/s^2}{1/(\sigma^2\Delta t)+1/s^2},\dfrac1{1/(\sigma^2\Delta t)+1/s^2}\right)$.

\subsubsection*{Conditional parameter updates}

Given the latent $\Delta N_i,Y_i$, use normal updates for $\mu,k$, inverse-gamma updates for $\sigma^2,s^2$, and a beta update for the jump probability. In particular,

$\sigma^2\mid\cdots\sim IG\left(\alpha+\dfrac n2,\beta+\dfrac{\sum_i(x_i-\mu\Delta t-Y_i\Delta N_i)^2}{2\Delta t}\right)$,

$\lambda\Delta t\mid\cdots\sim Be(a+N,b+n-N),\qquad N=\sum_i\Delta N_i$,

$k\mid\cdots\sim N\left(\dfrac{\tau_Y^2\bar Y+m_Ys^2/N}{\tau_Y^2+s^2/N},\dfrac{\tau_Y^2s^2}{N\tau_Y^2+s^2}\right)$,

$s^2\mid\cdots\sim IG\left(\alpha_Y+\dfrac N2,\beta_Y+\dfrac12\sum_{j=1}^N(Y_j-k)^2\right)$.

One Gibbs sweep updates $\mu,\sigma^2,\lambda,k,s^2$, every $\Delta N_i$, and then $Y_i$ where $\Delta N_i=1$. The Dow Jones example finds about $1.75$--$5.5$ jumps per year and large jump variances, so jump risk is material.

\subsection{Metropolis--Hastings}

\subsubsection*{Detailed balance}

For transition probabilities $p_{ij}$, stationarity requires $\pi(j)=\sum_i\pi(i)p_{ij}$. Detailed balance,
$\pi(i)p_{ij}=\pi(j)p_{ji}$,
is sufficient. Starting with proposal $q_{ij}$, accept $i\to j$ with
$\alpha_{ij}=\min\left\{1,\dfrac{\pi(j)q_{ji}}{\pi(i)q_{ij}}\right\}$.

Then $p_{ij}=q_{ij}\alpha_{ij}$ for $i\ne j$, with leftover probability assigned to staying at $i$; the resulting chain is reversible with stationary distribution $\pi$.

\subsubsection*{Continuous-state algorithm}

At state $\theta_{k-1}$, propose $X_k\sim q(\cdot\mid\theta_{k-1})$ and compute

$\alpha=\min\left\{1,\dfrac{\pi(X_k)q(\theta_{k-1}\mid X_k)}{\pi(\theta_{k-1})q(X_k\mid\theta_{k-1})}\right\}$.

Accept if $U\le\alpha$; otherwise set $\theta_k=\theta_{k-1}$. Symmetric proposals cancel the $q$ ratio, and unknown target normalizing constants cancel. Gibbs sampling is MH with full-conditional proposals and acceptance probability $1$.

\subsection{Stochastic Volatility And MH}

Let $y_t=e^{h_t/2}\epsilon_t$ and $h_{t+1}=\mu+\tau\eta_t$, with independent standard normal shocks. For priors $\mu\sim N(\alpha_\mu,\beta_\mu^2)$ and $\tau^2\sim IG(\alpha_\tau,\beta_\tau)$,

$\mu\mid\tau^2,h\sim N(\hat\alpha_\mu,\hat\beta_\mu^2),\quad \hat\alpha_\mu=\dfrac{\bar h\beta_\mu^2+\alpha_\mu\tau^2/n}{\beta_\mu^2+\tau^2/n},\quad \hat\beta_\mu^2=\dfrac{\beta_\mu^2\tau^2}{n\beta_\mu^2+\tau^2}$,

$\tau^2\mid\mu,h\sim IG\left(\alpha_\tau+\dfrac n2,\beta_\tau+\dfrac12\sum_{t=1}^n(h_t-\mu)^2\right)$.

The latent-state conditional is nonstandard:

$p(h_t\mid\mu,\tau,y,h_{-t})\propto e^{-h_t/2}\exp\left(-\dfrac{y_t^2}{2e^{h_t}}\right)\exp\left(-\dfrac{(h_t-\mu)^2}{2\tau^2}\right)$.

Use an independence proposal $x_t\sim q$, e.g. normal, and accept with

$\min\left\{1,\dfrac{p(x_t\mid\mu,\tau,y,h_{-t})q(h_t^{(i-1)})}{p(h_t^{(i-1)}\mid\mu,\tau,y,h_{-t})q(x_t)}\right\}$.

This is Metropolis-within-Gibbs: use Gibbs for tractable conditionals and MH for nonstandard latent-volatility conditionals.
